\section{Lecture 22}

\subsection{Multicollinearity}

\begin{definition}
\textbf{Multicollinearity}: This occurs when one (or more) of the predictors is nearly a linear combination of some of the others.

This leads to higher variance of the estimates of the regression coefficients. The coefficients may have large standard errors and become statistically insignificant, even when the overall F-statistic is significant.
\end{definition}

\subsection{Variance Inflation Factor (VIF)}

\begin{definition}
Consider the regression of the variable $X_j$ on the remaining predictors $\{X_k : k \neq j\}$.  
Let $R_j^2$ denote the coefficient of determination in this regression.

The \textbf{variance inflation factor} of $X_j$ is defined as
\[
    \mathrm{VIF}_j = \frac{1}{1 - R_j^2}.
\]
The VIF measures how much the variance of the LS estimate is inflated due to collinearity.
\end{definition}

\subsection{Interpretation of VIF}

Consider the simple (univariate) regression of $Y$ on $X_j$:
\[
    Y_i = \alpha_0 + \alpha_j X_{ji} + \varepsilon_i.
\]

Then
\[
    \Var(\hat{\alpha}_j) = \frac{\sigma^2}{s_j^2},
\]
where
\[
    s_j^2 = \sum_{i=1}^n (X_{ji} - \bar{X}_j)^2.
\]

\begin{result}
In the full regression model
\[
    Y = X \beta + \varepsilon,
\]
using the formula for the inverse of a partitioned matrix, one can show
\[
    \Var(\hat{\beta}_j) = \frac{\sigma^2}{1 - R_j^2} \cdot \frac{1}{s_j^2}.
\]
Thus the variance of $\hat{\beta}_j$ is inflated relative to the univariate case by the factor
\[
    \mathrm{VIF}_j = \frac{1}{1 - R_j^2}.
\]
\end{result}

\subsection{Methods to Handle Multicollinearity}

\begin{itemize}
    \item Remove some variables (to reduce collinearity in the remaining ones).
    \item Use shrinkage methods such as ridge regression.
\end{itemize}

\subsection{Ridge Regression}

Consider the multiple linear regression model
\[
    Y_i = \beta_0 + \beta_1 X_{1i} + \cdots + \beta_p X_{pi} + \varepsilon_i,
\]
for $1 \leq i \leq n$, where $p$ may be larger than $n$.

Given a regularization parameter $\lambda > 0$, the ridge regression estimator is
\[
    \hat{\beta}(\lambda)
    = \arg\min_{\beta}
    \left\{
        \sum_{i=1}^n \left( y_i - \beta_0 - \sum_{j=1}^p \beta_j X_{ji} \right)^2
        +
        \lambda \sum_{j=1}^p \beta_j^2
    \right\}.
\]

\paragraph{Remarks.}
\begin{enumerate}
    \item Larger $\lambda$ means more shrinkage toward zero.
    \item If $\lambda = 0$, we recover the least squares estimator.
    \item The intercept term is \textbf{not} penalized.
\end{enumerate}

\subsection{Centering and Scaling}

Typically the predictors are scaled to mean $0$ and variance $1$:
\[
    \bar{X}_j = \frac{1}{n} \sum_{i=1}^n X_{ji},
    \qquad
    s_j^2 = \frac{1}{n-1} \sum_{i=1}^n (X_{ji} - \bar{X}_j)^2,
\]
and the standardized version is
\[
    \tilde{X}_{ji} = \frac{X_{ji} - \bar{X}_j}{s_j}.
\]

This removes the intercept and makes ridge coefficients equivariant under scaling transformations
\[
    X_{ji} \mapsto a + b X_{ji}.
\]

\subsection{Constrained Form of Ridge Regression}

The equivalent constrained optimization problem is
\[
    \hat{\beta}(\lambda)
    =
    \arg\min_{\beta}
    \sum_{i=1}^n
        \left( y_i - \beta_0 - \sum_{j=1}^p \beta_j X_{ji} \right)^2
    \quad \text{subject to } \sum_{j=1}^p \beta_j^2 \le t.
\]

\subsection{Matrix Form}

Assume $Y$ is centered and $X$ is scaled so that $X$ has $p$ columns. Then the ridge objective is
\[
    RSS(\lambda)
    =
    \| Y - X \beta \|^2 + \lambda \|\beta\|^2.
\]

Then
\[
    \hat{\beta}(\lambda)
    = \arg\min_{\beta} RSS(\lambda).
\]

\begin{result}
The ridge regression estimator is
\[
    \hat{\beta}(\lambda)
    = (X^{\top} X + \lambda I_p)^{-1} X^{\top} Y,
\]
where $I_p$ is the $p \times p$ identity matrix.

\begin{enumerate}
    \item $\hat{\beta}(\lambda)$ is a linear function of $Y$.
    \item $X^{\top} X + \lambda I_p$ is invertible even when $X^{\top} X$ is not.
\end{enumerate}
\end{result}

\begin{proof}
Taking the gradient of the objective:
\[
    \nabla_{\beta} RSS(\lambda)
    = -2 X^{\top} Y + 2 X^{\top} X \beta + 2 \lambda \beta = 0.
\]
Thus
\[
    (X^{\top} X + \lambda I_p)\hat{\beta}(\lambda) = X^{\top} Y,
\]
and hence
\[
    \hat{\beta}(\lambda)
    = (X^{\top} X + \lambda I_p)^{-1} X^{\top} Y.
\]
\end{proof}

\newpage